Domain-specialized company memory and action recommendation pose different evaluation problems: a system must first provide sufficient current context, then use it to make an appropriate decision. We examine this boundary on a controlled B2B marketing research-action task with verified fictional records. The original frozen study uses 12 company worlds, 96 queries, six evidence-selection adapters and one shared quantized Qwen reader. At a 1,024-token allowance, catalog-conditioned coverage yields 53/96 correct supported decisions versus 40/96 for prompt-specialized hybrid retrieval: a 13.5-point difference with a world-bootstrap interval of [0.0, 28.1], which does not establish clear superiority. Complete support is available in all 48 answerable cases, yet the reader produces only 21 supported recommendations. An exploratory extension executes the actual GBrain hybrid-search component on the same records. Applying the same prerequisite-selection module to its results closes the support gap; larger context allowances also let generic controls recover complete support. An exact SQL reference solves the known formal rules without a model. These findings favor separating context quality from recommendation quality and evaluating specialization as a composable layer. They do not establish a need to replace general company memory, strategic usefulness, raw-document ingestion quality or commercial advantage.
